\chapter{Introduction}
\label{chap:intro}

Global energy investment exceeded 3 trillion USD in 2024~\cite{iea_investment_2024}, and the ability to forecast the performance of hydrocarbon assets remains a key factor of project viability. Numerical reservoir simulation serves as the foundational instrument for such forecasting. By constructing \enquote{digital twins} of the subsurface, engineers predict production dynamics and optimize recovery strategies over time horizons of decades. However, the application of reservoir simulation is inherently constrained by the complexity of the geological systems it attempts to represent. As noted by Saleri and Toronyi~\cite{spe_recovery_economics}, simulation should not be viewed as a deterministic calculator but rather as a probabilistic tool subject to substantial uncertainties and non-uniqueness. The reliability of these models depends on their ability to capture the spatial distribution of rock properties and their impact on fluid flow~\cite{ringrose_value_2013}. Ringrose and Bentley emphasize that the primary purpose of a reservoir model is to provide predictive capacity by placing physical properties correctly in three-dimensional space, rather than merely generating visually appealing graphics.

Reservoir simulation, however, faces an intrinsic contradiction between predictive accuracy and computational cost. High-fidelity models, which are required to resolve fine-scale geological heterogeneities and complex multiphase flow phenomena, demand excessive computational resources. Coarse-grid models, although computationally efficient, often fail to capture essential fluid dynamics and produce systematic forecasting errors. The mismatch between the operational need for rapid, reliable forecasts and the limitations of classical simulators creates a substantial bottleneck in effective reservoir management.

To mitigate these uncertainties and maximize asset value, the industry has increasingly adopted \enquote{closed-loop reservoir management}~\cite{jansen_closed_loop_2009}. This paradigm shifts reservoir management from a periodic activity to a near-continuous process, requiring the rapid evaluation of numerous scenarios to optimize well controls and update geological models in real-time. The practical implementation of this approach, however, varies significantly across the industry. One prevalent strategy, actively pursued by Russian operators such as Gazprom Neft, focuses on integrated digital platforms. As described by Kusimova et al.~\cite{gazpromneft_ml_conference}, machine learning in this context is primarily employed to automate peripheral processes surrounding the hydrodynamic simulator, such as interpreting seismic data or predicting well failures. While this enhances overall workflow efficiency, the core prediction engine remains a classical numerical simulator, and the cost of each individual high-fidelity forecast stays unchanged.

A contrasting philosophy is observed among Western majors, such as Equinor, which address geological uncertainty through ensemble modeling. Rather than seeking a single deterministic forecast, this strategy assesses a probabilistic range of outcomes using data assimilation methods like Ensemble Smoother with Multiple Data Assimilation (ES-MDA)~\cite{equinor_es_mda}. Technologically, this approach necessitates launching thousands of parallel simulations to generate probabilistic production profiles. The requirement to run such massive ensembles, while effective for uncertainty quantification, sits at the limit of current computational capabilities.

The computational cost of such ensemble-based workflows is prohibitive when applied to high-fidelity geological models. Hayder et al.~\cite{stanford_hpc_simulation} report that running models with billions of cells requires massive high-performance computing (HPC) infrastructure and can take weeks to complete a single realization. To make ensemble calculations feasible, practitioners typically upscale geological models, increasing grid block sizes from the required fine scale (e.g., $10\times10$ meters) to coarse scales (e.g., $100\times100$ meters)~\cite{slb_delfi_ensemble}. This upscaling process introduces a trade-off: while it reduces simulation runtime, it leads to the loss of information regarding fine-scale geological heterogeneities, such as high-permeability channels and flow barriers. Reviews of uncertainty quantification methods~\cite{uq_reservoir_review} attribute this loss to the inability of coarse models to resolve such features, resulting in numerical diffusion and a failure to reproduce complex, unstable flow phenomena like viscous fingering.

Although reduced-order modeling (ROM) techniques, such as Trajectory Piecewise Linearization (TPWL), have been proposed to accelerate these calculations by factors of 100 or more~\cite{rom_review_spe}, they often require computationally expensive offline training and may lack robustness when applied to highly nonlinear multiphase flow problems outside the training regime. Grid coarsening therefore remains the dominant strategy for practical ensemble modeling. However, recent studies by Pola et al.~\cite{geological_heterogeneity_impact} demonstrate that this strategy is inadequate for complex recovery processes. The numerical dispersion induced by grid coarsening can introduce errors comparable in magnitude to those caused by an incorrectly configured physical model. Coarse grids fail to resolve micro- and meso-scale heterogeneities that govern the propagation of fluid fronts. From a mathematical perspective, the dynamics of multiphase filtration in fractured-porous media are described by nonlinear parabolic partial differential equations, structurally analogous to the nonlinear diffusion problems investigated in recent super-resolution studies~\cite{konyukhov_article}. The resulting forecasts are systematically biased and cannot be corrected by simple parameter adjustment. Traditional deterministic reconstruction techniques, such as bilinear and bicubic interpolation, do not address this issue adequately either: recent benchmarks~\cite{konyukhov_article} show that these methods effectively act as low-pass filters, smoothing out sharp discontinuities and failing to recover the high-frequency geometric details of flow boundaries lost during coarsening. The field therefore lacks computational tools capable of recovering high-fidelity, physically consistent flow details from coarse-grid simulations without incurring the cost of fine-grid modeling. Deep learning, and super-resolution algorithms in particular, offers a candidate methodology for this problem by learning the mapping between large-scale dynamics and fine-scale physics.

The primary goal of this thesis is to develop and investigate a neural network software complex for the accelerated reconstruction of transient high-fidelity pressure and saturation fields in fractured-porous media. The proposed solution combines residual super-resolution backbones with adversarial training, physics-informed regularization, and conditioning on the dynamic reservoir state. Spatial super-resolution is applied independently at each simulation time step, and the sequence of reconstructed states reproduces the temporal evolution of the multiphase flow.

The experimental framework of this study relies on a synthetic dataset of two-dimensional simulations in fractured-porous reservoirs, in which macroscopic flow occurs through a highly conductive fracture network while porous matrix blocks act as fluid sources. To avoid non-physical downsampling artifacts and to capture numerical diffusion over time accurately, the paired coarse-grid low-resolution states and fine-grid high-resolution states are generated by independent numerical solvers rather than by direct subsampling. The continuous temporal evolution of the reservoir is discretized into sequential snapshots. Each temporal state is then formulated as a three-channel physical tensor in which the first channel represents the pressure field, the second channel represents the water saturation in the fracture network, and the third channel represents the water saturation in the porous matrix blocks. This dual-saturation formulation explicitly reflects the dual-porosity nature of the modeled medium and preserves the distinction between fluid storage in matrix blocks and fluid transport along the fractures.

To evaluate model performance beyond purely visual assessment, the present work employs a suite of complementary metrics. The Peak Signal-to-Noise Ratio (PSNR) is used to measure point-wise mathematical accuracy across the computational grid, while the Structural Similarity Index (SSIM) evaluates the preservation of macroscopic flow geometries and spatial connectivity within the reservoir. The Error Relative Global Adimensional Synthesis (ERGAS) together with the Spectral Angle Mapper (SAM) quantify the radiometric quality of the multi-channel reconstruction and the preservation of inter-channel physical relationships between the pressure and saturation fields. A dedicated gradient-based criterion further assesses the recovery of sharp displacement fronts, which represent the high-frequency spatial features typically lost to numerical diffusion during grid coarsening. Finally, to verify the physical plausibility of the reconstructed states, a physics-informed consistency criterion penalizes deviations from the diffusive nature of the pressure field and from the geometric coherence of saturation fronts across the dual-porosity system. Based on this evaluation framework, the research addresses the following Research Questions (RQs):
\begin{itemize}
    \item RQ1: How does the choice of optimization objective, specifically pixel-wise regression versus adversarial training augmented with a perceptual criterion, quantitatively affect the trade-off between point-wise accuracy and structural-spectral fidelity when reconstructing transient pressure and saturation fields in fractured-porous reservoirs?
    \item RQ2: To what extent does the integration of physics-informed consistency criteria, namely the smoothness regularization of the pressure field and the gradient-magnitude alignment of saturation fronts, reduce non-physical artifacts in adversarial super-resolution outputs compared to unconstrained generative training?
    \item RQ3: To what extent does conditioning the super-resolution generator on dynamic simulation scalars, such as injection pressure, cumulative water cut, and well operational regime, improve reconstruction fidelity compared to an unconditional baseline operating solely on the coarse-grid state tensor?
\end{itemize}

Based on these questions, the following hypotheses are formulated:

Hypothesis 1 (Structural Sharpness and Point-wise Fidelity Trade-off). The application of an adversarial network framework constrained by perceptual loss recovers the sharp fluid displacement fronts artificially smoothed by conventional regression models. This architectural shift improves structural and spectral fidelity at the cost of lower point-wise accuracy, which reproduces the perception-distortion trade-off known from the broader super-resolution literature.
\begin{itemize}
    \item Independent Variable: Optimization objective, specifically point-wise loss versus adversarial perceptual loss.
    \item Dependent Variable: The measured trade-off between structural or spectral evaluation scores and point-wise error metrics.
    \item Population: Spatial distributions of multiphase flow state variables within the studied computational domain.
    \item Relation: Adversarial training with a perceptual criterion biases the network towards reconstructions that preserve sharp displacement fronts and resolve numerical diffusion, at the price of higher per-pixel error relative to a pixel-loss baseline.
\end{itemize}

Hypothesis 2 (Physics-Informed Consistency). The integration of a physics-informed regularization term, which combines a smoothness penalty on the pressure field with a gradient-magnitude alignment constraint on the saturation channels, into the adversarial framework reduces non-physical artifacts compared to unconstrained generative training. This physical regularization prevents the synthesis of spatially inconsistent fluid states in fractured-porous media without sacrificing the structural sharpness gained from adversarial training.
\begin{itemize}
    \item Independent Variable: The presence and weight of the physics-informed consistency term within the adversarial optimization objective.
    \item Dependent Variable: The magnitude of non-physical artifacts in the reconstructed fields, quantified by gradient-based and maximum-error metrics.
    \item Population: Spatial distributions of multiphase flow state variables within the studied computational domain.
    \item Relation: The combined smoothness and gradient-alignment terms act as a regularizer that pushes the network towards pressure fields consistent with the diffusion equation and towards saturation fields with well-defined fronts.
\end{itemize}

Hypothesis 3 (Conditioning on Dynamic Reservoir State). Conditioning the super-resolution generator on a vector of dynamic simulation scalars, propagated into the residual feature stream through a feature-wise linear modulation mechanism, yields a measurable improvement in reconstruction fidelity compared to an unconditional baseline. The conditioning signal supplies information about the current operational regime of the reservoir that the coarse-grid state alone does not unambiguously determine, and thus reduces the ill-posedness of the upscaling problem.
\begin{itemize}
    \item Independent Variable: The presence and composition of the dynamic conditioning vector supplied to the generator.
    \item Dependent Variable: Reconstruction fidelity quantified by point-wise, structural, and spectral evaluation metrics.
    \item Population: Spatial distributions of multiphase flow state variables within the studied computational domain.
    \item Relation: Feature-wise linear modulation conditioned on dynamic scalars supplies the generator with information about the current well regime, which is otherwise unavailable from the coarse-grid input and may correspond to several distinct high-resolution states.
\end{itemize}

To achieve the goal and test the formulated hypotheses, the following objectives have been defined:
\begin{enumerate}
    \item Conduct a domain analysis to identify the limitations of coarse-grid numerical modeling in fractured-porous media and review existing neural network approaches to spatial super-resolution.
    \item Implement and evaluate three residual super-resolution architectures, namely MSRResNet, MDSR, and RFDN, optimized with a pixel-wise loss to establish point-wise regression baselines for the three-channel state tensors.
    \item Implement an adversarial training framework augmented with a perceptual criterion based on gradient and spectral consistency, in order to recover high-frequency structural details and evaluate the perception-distortion trade-off on physical flow fields.
    \item Develop a physics-informed regularization module that combines a pressure smoothness penalty with a saturation gradient-alignment constraint, and integrate it into the adversarial objective function as an auxiliary loss term.
    \item Extend the MDSR baseline with a feature-wise linear modulation mechanism conditioned on dynamic simulation scalars, and conduct an ablation study over alternative conditioning vector configurations to assess the contribution of the conditioning signal.
    \item Perform computational experiments to compare the proposed methodology against pixel-loss and unconditional baselines using an evaluation framework that includes point-wise accuracy metrics (PSNR, MAE), structural and spectral metrics (SSIM, ERGAS, SAM, gradient-based front consistency), and physics-informed consistency metrics.
    \item Develop a prototype desktop application based on the PySide6 framework with integration to the simulation backend via a gRPC interface, in order to demonstrate the practical inference of the best-performing model and enable side-by-side comparison of classical interpolation, neural super-resolution, and ground-truth fine-grid solutions.
\end{enumerate}

The practical significance of this work consists in reducing the computational cost of high-fidelity reservoir forecasting through partial replacement of fine-grid simulation runs by neural super-resolution applied to coarse-grid solutions. The proposed methodology enables larger ensemble sizes in uncertainty quantification workflows without proportional growth in high-performance computing requirements, which makes probabilistic decision-making more accessible in field development planning. The accompanying desktop application integrates with the existing simulation infrastructure via a gRPC interface and allows engineers to compare classical interpolation, neural super-resolution, and ground-truth fine-grid solutions within a unified visual environment.

The thesis consists of an introduction, six chapters, a conclusion, a bibliography, and appendices. Chapter 2 presents a detailed literature review, covering reservoir geology, numerical simulation, and neural network methods. Chapter 3 details the research methodology, experimental plan, and architecture descriptions. Chapter 4 describes the practical implementation of the software complex, including data generation and optimization. Chapter 5 presents the results and discussion, analyzing model performance and interpretability. Finally, Chapter 6 summarizes the main findings and outlines directions for future research.
